% SPDX-License-Identifier: Apache-2.0
\section{An ALU as a Stream, Out of Order}
\label{sec:x-stream-alu}

An ALU that takes its work as a stream: a request is a tag, an
operation and two operands on one channel, and a result is the tag
and the value on another. The four operations are four units of four
depths. \code{not} takes one cycle, \code{and} two, \code{or} three
and \code{add} four, a nibble per stage with the carry passed from
one stage to the next. A result leaves as soon as its unit is done,
so results come out in the order the units finish and not the order
the requests came, and the tag says which request each one answers.

One result can leave per cycle. So a request is taken only when the
cycle its unit would finish in is free: \code{slots} holds a bit per
cycle ahead, a request claims the bit of its latency, and the
register shifts down every cycle. A request whose bit is claimed
waits at the head of the channel, its \code{ready} low, until the
bit is free. The sink's room moves every stage at once; when the
sink holds off, every stage holds and nothing is taken. The unit is
lowered, and its netlist is checked against this run under both
simulators.

The run offers ten requests. The first is an add, and the \code{not}
behind it comes out two cycles before it; the \code{and} behind that
would finish in the add's cycle, so it waits one cycle at the head.
The sink holds off for three cycles in the middle, and the whole
pipeline holds with it.

\begin{figure}[htbp]
\centering
\input{stream_alu_timing.tex}
\caption{The ALU as a stream: the requests offered and taken, the
claimed slots, the four units' last stages, and the results as they
leave, out of order.}
\label{fig:x-stream-alu}
\end{figure}

\lstinputlisting[caption={\code{ex\_stream\_alu.rs}. Run with \code{bazel run //lib/examples:ex\_stream\_alu}.},label={lst:x-stream-alu}]{lib/examples/ex_stream_alu.rs}

\lstinputlisting[style=ascii,caption={What \code{ex\_stream\_alu} printed when the build ran it: what the source sent, the slots, what the first stage of a unit holds, and the result taken, per cycle; then the lowering.},label={lst:x-stream-alu-out}]{out_stream_alu.txt}
